\section{Endomorphism algebras and Weil type}\label{app:albert}

\Cref{def:weiltype} asks for three things: an even dimension, an imaginary
quadratic field acting with balanced eigenvalue multiplicities, and
compatibility with the polarisation. This appendix places these conditions
inside Albert's classification of endomorphism algebras. Two statements used
in the body come out of it. The first is that the compatibility condition is
automatic once the field is stable under the Rosati involution, so that it
reduces to that stability (\Cref{lem:rosatiauto}). The second is that the
balanced condition is not a normalisation but exactly the condition under
which the Weil line carries a Hodge class (\Cref{prop:signatureforced}); this
is why the subject concerns balanced signatures alone.

\subsection{Positive involutions}

\begin{definition}\label{def:positive}
Let $D$ be a finite dimensional semisimple $\QQ$-algebra. An involution
$\dagger$ of $D$, that is a $\QQ$-linear anti-automorphism with
$\dagger^{2}=\id$, is \emph{positive} if
\[
  \Tr_{D/\QQ}\bigl(x\,x^{\dagger}\bigr)>0
  \qquad\text{for every }x\in D,\ x\ne0,
\]
where $\Tr_{D/\QQ}$ is the trace of left multiplication on $D$ as a
$\QQ$-vector space.
\end{definition}

For an abelian variety $A$ with a polarisation $\eta$, the algebra
$D=\End^{0}(A)=\End(A)\otimes\QQ$ is semisimple and the Rosati involution
$\dagger$ attached to $\eta$ is positive \cite[Ch.~5]{BL04}. Everything below
is a statement about the pair $(D,\dagger)$.

\begin{lemma}[Totally real centre]\label{lem:totallyreal}
Let $(D,\dagger)$ be as above with $D$ a division algebra, let $F$ be the
centre of $D$ and $F_{0}=\{z\in F:z^{\dagger}=z\}$. Then $F_{0}$ is a totally
real number field.
\end{lemma}

\begin{proof}
For $z\in F_{0}$ the positivity condition reads
$\Tr_{D/\QQ}(z^{2})>0$, and $\Tr_{D/\QQ}$ restricted to $F_{0}$ is
$[D:F_{0}]\cdot\Tr_{F_{0}/\QQ}$, so $\Tr_{F_{0}/\QQ}(z^{2})>0$ for every
nonzero $z\in F_{0}$. Suppose $F_{0}$ had a complex place $\sigma$. The
embedding of $F_{0}$ into the product of its archimedean completions has dense
image, so there is $z\in F_{0}$ with $\sigma(z)$ within $\tfrac1{10}$ of $i$
and $|\rho(z)|<\varepsilon$ for every other archimedean place $\rho$. Then
$\Tr_{F_{0}/\QQ}(z^{2})=\sum_{\rho}\rho(z)^{2}$, where a complex place
contributes twice the real part of $\rho(z)^{2}$. The real part of
$\sigma(z)^{2}$ is less than $-1+\tfrac{2}{10}+\tfrac{1}{100}<-\tfrac12$, so
the term from $\sigma$ is less than $-1$, and the remaining terms are
bounded in absolute value by $2[F_{0}:\QQ]\varepsilon^{2}$, so the total is
negative for $\varepsilon$ small. This contradicts positivity, so every place
of $F_{0}$ is real.
\end{proof}

\begin{theorem}[Albert \cite{Alb34,Alb35a,Alb35b}]\label{thm:albert}%
\footnote{The numbering of the four types is Albert's and has stuck, although
it follows no natural order: types I to III have a totally real centre and
type IV a CM one, so the natural division is two against two rather than three
against one. Shimura's normalisation of the same classification, in terms of
the associated symmetric domains, is the one that makes the moduli side
transparent.}
Let $(D,\dagger)$ be a division algebra over $\QQ$ with a positive involution,
$F$ its centre, $F_{0}$ the subfield fixed by $\dagger$, $e_{0}=[F_{0}:\QQ]$
and $[D:F]=d^{2}$. Then $F_{0}$ is totally real, $[F:F_{0}]\in\{1,2\}$, and
exactly one of the following holds.
\begin{enumerate}[label=\textup{Type \Roman*.},leftmargin=4.2em,itemsep=2pt]
\item $D=F=F_{0}$ is a totally real field and $\dagger$ is the identity.
\item $F=F_{0}$ is totally real and $D$ is a quaternion algebra over $F$ with
  $D\otimes_{F,\rho}\RR\cong M_{2}(\RR)$ at every real place $\rho$;
  $\dagger$ is of orthogonal type.
\item $F=F_{0}$ is totally real and $D$ is a quaternion algebra over $F$ with
  $D\otimes_{F,\rho}\RR\cong\mathbb{H}$ at every real place $\rho$; $\dagger$
  is the canonical involution.
\item $F$ is a CM field with maximal totally real subfield $F_{0}$, and
  $\dagger$ induces complex conjugation on $F$.
\end{enumerate}
If moreover $D=\End^{0}(A)$ for an abelian variety $A$ of dimension $g$ with
$\dagger$ the Rosati involution of a polarisation, then
\[
  \text{\rm I: } e_{0}\mid g,
  \qquad
  \text{\rm II: } 2e_{0}\mid g,
  \qquad
  \text{\rm III: } 2e_{0}\mid g,
  \qquad
  \text{\rm IV: } e_{0}d^{2}\mid g .
\]
\end{theorem}

We quote \Cref{thm:albert} as a classification. \Cref{lem:totallyreal}
proves the part of it that we use directly, and the four types with their
divisibility conditions are in \cite[Ch.~9]{BL04}. \Cref{fig:albert} draws the
signatures of an imaginary quadratic action and the dimensions of their
families.

\begin{figure}[htbp]
\centering
  \includegraphics[width=0.92\textwidth]{fig_signature}
\caption{The signatures $(p,q)$ of an imaginary quadratic action, with
$p+q=\dim A$, drawn on the surface of height $pq$, which is the dimension of
the family of abelian varieties carrying that action, the unitary Shimura
variety of signature $(p,q)$. Each curve on the surface collects the
signatures available in one dimension, and on it the diagonal point is the
highest. The divisibility conditions of \Cref{thm:albert} put no
constraint on which of them occur when $\End^{0}(A)=K$, by
\Cref{prop:whichtypes}(i). By \Cref{prop:signatureforced} the Weil line
consists of Hodge classes exactly at the filled points on the ridge
$p=q$, and at every open point it is a rational sub-Hodge structure of type
$\{(p,2n-p),(2n-p,p)\}$ carrying no Hodge class at all. The diagonal is the
locus that \Cref{cor:weilalgebraic} shows to be an algebraic subvariety of the
moduli.}
\label{fig:albert}
\end{figure}

\subsection{Compatibility with the polarisation}

\begin{lemma}[Rosati on an imaginary quadratic subfield]\label{lem:rosatiauto}
Let $A$ be an abelian variety with a polarisation $\eta$ and Rosati involution
$\dagger$, and let $\iota\colon K\hookrightarrow\End^{0}(A)$ be an embedding of
an imaginary quadratic field such that $\iota(K)$ is stable under $\dagger$.
Then
\[
  \iota(\tau)^{\dagger}=\iota(\bbar{\tau})
  \qquad\text{for every }\tau\in K .
\]
That is, condition \textup{(iii)} of \Cref{def:weiltype} holds automatically,
and Weil type is the conjunction of \textup{(i)}, \textup{(ii)} and the
$\dagger$-stability of $\iota(K)$.
\end{lemma}

\begin{proof}
Put $D=\End^{0}(A)$. By stability, $\dagger$ restricts to an involutive
anti-automorphism of the field $\iota(K)\cong K$. An anti-automorphism of a
commutative ring is an automorphism, and $\Gal(K/\QQ)$ has order two, so the
restriction is either the identity or complex conjugation. Suppose it were
the identity, and put $x=\iota(\sqrt{-d})$. Then
$x\,x^{\dagger}=x^{2}=\iota(-d)=-d\cdot 1$, so
\[
  \Tr_{D/\QQ}\bigl(x\,x^{\dagger}\bigr)=-d\cdot\dim_{\QQ}D<0,
\]
contradicting the positivity of $\dagger$. Hence $\dagger$ restricts to complex
conjugation.
\end{proof}

\begin{remark}\label{rem:whystable}
Stability is a real hypothesis: a subfield of $\End^{0}(A)$ need not be
$\dagger$-stable. What \Cref{lem:rosatiauto} removes is the second half of the
condition, the identification of $\dagger|_{K}$ with a specific automorphism.
In particular, in the model of \Cref{setup:model} the polarisation need not be
chosen to match the field, since any $\dagger$-stable embedding does the same
work.
\end{remark}

\subsection{The balanced signature}

\begin{definition}[Signature]\label{def:signature}
Let $\iota\colon K\hookrightarrow\End^{0}(A)$ be as above, $g=\dim A$. The
operator $\iota(\sqrt{-d})$ acts on the tangent space $T_{0}A\cong\CC^{g}$ with
the two eigenvalues $\pm i\sqrt{d}$; write $p$ and $q=g-p$ for their
multiplicities. The pair $(p,q)$ is the \emph{signature} of $(A,\iota)$.
\end{definition}

The signature is an isogeny invariant, since an isogeny induces an isomorphism
on tangent spaces commuting with the action. It is locally constant in
families: the two eigenspaces are the kernels of
$\iota(\sqrt{-d})\mp i\sqrt{d}$ on the relative tangent bundle, so their
dimensions are upper semicontinuous, and they sum to $g$. By
duality $\iota(\tau)^{*}$ acts on $H^{1,0}(A)$ with the same multiplicities, so
$\dim V^{1,0}_{+}=p$ and $\dim V^{1,0}_{-}=q$ in the notation of
\Cref{sec:weiltype}.

\begin{proposition}[The balanced signature is forced]\label{prop:signatureforced}
Let $A$ be an abelian variety of dimension $2n$ with a $\dagger$-stable
embedding $\iota\colon K\hookrightarrow\End^{0}(A)$ of signature $(p,q)$,
$p+q=2n$, and let
\[
  \HW(A)=\Bigl(\textstyle\bigwedge^{2n}_{\CC}V_{+}
   \oplus\bigwedge^{2n}_{\CC}V_{-}\Bigr)\cap H^{2n}(A,\QQ)
\]
as in \Cref{def:weilline}. Then $\HW(A)$ is a rational sub-Hodge structure of
$H^{2n}(A,\QQ)$ of dimension two and of Hodge type
$\{(p,2n-p),\,(2n-p,p)\}$, and the following are equivalent:
\begin{enumerate}[label=\textup{(\roman*)},nosep]
\item $p=q=n$, that is $(A,\iota,\eta)$ is of $(K,-1,n)$-Weil type;
\item $\HW(A)$ contains a nonzero Hodge class;
\item $\HW(A)$ consists of Hodge classes.
\end{enumerate}
\end{proposition}

\begin{proof}
The space $\bigwedge^{2n}V_{+}$ is a line, and by \eqref{eq:bisplit} it equals
$\bigwedge^{p}V^{1,0}_{+}\otimes\bigwedge^{2n-p}V^{0,1}_{+}$, which is of pure
Hodge bidegree $(p,2n-p)$; the conjugate line $\bigwedge^{2n}V_{-}$ is of
bidegree $(2n-p,p)$. Their sum is stable under conjugation and defined over
$\QQ$ by the argument of \Cref{prop:weilline}(i), which does not use the
balanced condition, so $\HW(A)$ is a two-dimensional rational sub-Hodge
structure of the stated type.

A Hodge class is an element of bidegree $(n,n)$. A nonzero element of
$\HW(A)\otimes\CC$ has a nonzero component in one of the two lines, so it is of
bidegree $(n,n)$ only if $(p,2n-p)=(n,n)$ or $(2n-p,p)=(n,n)$, and either
equality says $p=n$. This proves that (ii) implies (i). If $p=n$ then both
lines are of bidegree $(n,n)$ and every element of $\HW(A)$ is a Hodge class,
so (i) implies (iii); and (iii) implies (ii) because $\dim\HW(A)=2>0$.
\end{proof}

\begin{remark}\label{rem:onlybalanced}
\Cref{prop:signatureforced} is the reason the literature speaks of abelian
varieties of Weil type rather than of imaginary quadratic multiplication in
general. For an unbalanced signature the space $\HW(A)$ is still a canonical
two-dimensional rational sub-Hodge structure, and it is still not in the
divisor subring, but it contains no Hodge class and therefore predicts
nothing. The exceptional Hodge classes appear exactly at balance, and the
locus of balance is the algebraic subvariety of \Cref{cor:weilalgebraic}.
\end{remark}

\subsection{Endomorphism algebras of Weil type}

\begin{proposition}[Weil type in the classification]\label{prop:whichtypes}
Let $A$ be an abelian variety of dimension $2n$ admitting a $\dagger$-stable
embedding of an imaginary quadratic field $K$ with balanced signature, and
suppose $\End^{0}(A)$ is a division algebra $D$ with centre $F$ as in
\Cref{thm:albert}. Then:
\begin{enumerate}[label=\textup{(\roman*)},nosep]
\item if $D=K$ then $D$ is of type IV with $e_{0}=1$ and $d=1$, and the
  divisibility condition of \Cref{thm:albert} is vacuous, so every even
  dimension occurs;
\item the classification does not exclude type III with $F_{0}=\QQ$: a
  totally definite quaternion algebra over $\QQ$ contains imaginary
  quadratic subfields on which the canonical involution is conjugation;
\item $D$ is not of type I, because a totally real field contains no
  imaginary quadratic subfield.
\end{enumerate}
\end{proposition}

\begin{proof}
(i) If $D=K$ then $F=K$ is a CM field with $F_{0}=\QQ$, so $e_{0}=1$; and
$[D:F]=1$ gives $d=1$. The involution is conjugation by
\Cref{lem:rosatiauto}, which is type IV. The condition $e_{0}d^{2}\mid g$ reads
$1\mid 2n$.

(ii) Let $D$ be a quaternion algebra over $\QQ$ ramified at the infinite place,
with standard generators $u,v$ satisfying $u^{2}=a<0$, $v^{2}=b<0$, $uv=-vu$,
and canonical involution $x\mapsto\bbar{x}$ fixing $\QQ$ and negating
$u,v,uv$. Then $\QQ(u)\cong\QQ(\sqrt{a})$ is imaginary quadratic, stable under
the involution, and the involution restricts to conjugation on it. Whether
such a $D$ occurs as $\End^{0}(A)$ together with a balanced signature
is a separate question about the moduli, but nothing in the classification
excludes it.

(iii) is immediate, since a totally real field has no element of negative
square.
\end{proof}

\begin{remark}[The generic case]\label{rem:genericcase}
Case (ii) of \Cref{prop:whichtypes} is a special position: there
$\End^{0}(A)$ is larger than $K$, the Hodge group is correspondingly smaller,
and $\Hdg^{n}(A)$ can be larger than the three dimensions of
\Cref{thm:hdgcount}. It is also a case in which the Weil classes can become
accessible: \Cref{rem:special} records that at such points the Weil line can
meet the divisor subring, which is the split situation. The hard case, and the
one addressed throughout \Cref{sec:divisor,sec:characters,sec:hyperplane}, is
case (i) with $\Hg(A)=\SU(V,H)$.
\end{remark}

\begin{remark}[Dimension four]\label{rem:dimfour}
For $n=1$ the Weil line lies in $H^{2}(A,\QQ)$ and consists of divisor
classes, algebraic by the Lefschetz theorem on $(1,1)$ classes, so it predicts
nothing new; this is the same boundary as in \Cref{prop:normpart}. Moonen and
Zarhin determined the Hodge classes on abelian varieties of dimension at most
five \cite{MZ99}, and on simple abelian fourfolds in \cite{MZ95}: through
dimension three the Hodge ring is generated by divisor classes, and the first
Hodge classes not of that form occur in dimension four, on abelian fourfolds of
Weil type. Thus $n=2$ is at once the first case of \Cref{def:weiltype} with
exceptional Hodge classes, the first case where \Cref{thm:hdgcount} gives
three rather than one, and the first case of the Hodge conjecture for
abelian varieties that needed more than divisor classes; Markman settled it
\cite[Corollary 1.6.1]{Mar25a}. The next case is $n=3$ with nontrivial
discriminant, by \Cref{prop:separate}(ii).
\end{remark}
